\subsection{Closing degree thirty-two and the first nine degrees}
\label{sec:k32-vanishing}

\begin{theorem}[K17-K32-VANISH]\label{thm:k32-vanish}
Under the stipulated arithmetic theorem of Alp\"oge,
\[
                         K_{32}(\mathbf Z)=0.
 \tag{K32.1}
\]
Together with the preceding results, this gives
\[
             K_{4k}(\mathbf Z)=0\qquad(1\leq k\leq9).
 \tag{K32.2}
\]
For every odd prime $\ell$, one has
$A_\ell^{[-16]}=0$ and $\mathcal S_\ell(16)=0$.
The full-domain Soul\'e character
\[
 \kappa_{17}:N_\ell\longrightarrow\mathbf Z_\ell(17),
 \qquad N_\ell=\operatorname{Gal}
 (\overline{\mathbf Q}/\mathbf Q(\mu_{\ell^\infty})),
\]
is surjective, as are all its reductions modulo $\ell^a$ for $a\geq1$.
Consequently the full-domain characters $\kappa_m$ are surjective
for every odd $m$ with $3\leq m\leq19$ and every odd $\ell$.
\end{theorem}

\begin{proof}
Corollary~\ref{cor:k32-oneprime} has already proved that
$K_{32}(\mathbf Z)$ is a finite $q$-primary group, where
$q=25078104767246023$. It therefore remains to prove the vanishing
of precisely this primary part.

Use $E=\mathbf Q(\sqrt{-7})$, $e=16$, $N=2$, and $\psi=1$.
Lemma~\ref{lem:minus7-source} verifies every condition of the actual
source datum and permits geometric Frobenius at the split prime $2$
in the intermediate direct unit criterion. Its use here does not
invoke the separate auxiliary-prime restrictions of the counting
theorem.

Theorem~\ref{thm:u3-minus7-unit} supplies an integral Hecke lattice
in the entire two-class compact space of highest weight $(15,0,-15)$
and proves that
\[
       \det\bigl((1+2^{16}+2^{32})I-2^{15}T_2\bigr)
\]
is a unit above $q$. Its full $171$-dimensional Schur comparison
has determinant
\[
                    13187546903051047\not\equiv0\pmod q.
\]
Lemma~\ref{lem:minus7-u3-coverage} places every actual admitted
irreducible rank-three system in that entire space with the required
Frobenius normalization, in either minuscule direction.

Proposition~\ref{prop:minus7-gl2} supplies the entire integral
weight-$33$ space $S_{33}(\Gamma_1(7))$ and proves
\[
 \det\bigl((1+2^{32})^3I-T_2^3\bigr)
       \equiv15133664954664570\not\equiv0\pmod q.
\]
Lemma~\ref{lem:minus7-gl2-conductor} shows that every actual
admitted rank-two system is covered after a global mod-$7$ character
twist; all untwisting values at $2$ lie in $\mu_3$.
Thus the cubic determinant covers every one of those possibilities,
including systems whose original untwisted conductor is $49$.

Both full-space unit hypotheses of Corollary~\ref{cor:minus7-unit}
are now satisfied. That corollary applies the actual finite-set
inequality of Lemma~\ref{lem:step17-unit} and gives source Selmer
length zero. For $\psi=1$, the source module is the rational-field
strict negative-character group, with coefficient extension if
needed. The bottom coefficient ring is $\mathbf Z_q$; a finite
coefficient DVR is faithfully flat over it, so zero length after
that extension implies zero at the bottom ring. The extra local
condition at $2$ is automatic by the unramified local calculation
in K17-FIELD-DESC. Equivalently, the precise identification already
proved in (K5) gives
\[
       \bigl(K_{32}(\mathbf Z)\{q\}\bigr)^\vee
                        =\mathcal S_q(16)=0.
\]
Hence the last primary part vanishes, proving (K32.1).
The known degrees four and eight, Theorem~\ref{thm:k12-vanish},
Theorem~\ref{thm:higher-k-vanish}, and
Theorem~\ref{thm:k28-k36-vanish} give (K32.2).

For every odd $\ell$, the exact class-group comparison (K2) now
gives $A_\ell^{[-16]}=0$, and the strict-duality comparison (K5)
gives $\mathcal S_\ell(16)=0$. The full-domain index formula (K18)
at $m=17$ gives
\[
 [\mathbf Z_\ell(17):\kappa_{17}(N_\ell)]
       =|K_{32}(\mathbf Z)\{\ell\}|=1.
\]
This proves surjectivity and its reductions. Applying the same
formula to every degree in (K32.2) gives the final weight range.
The normalization remains the cyclotomic product (K19), on the
stated full domain; no claim about unit values on a deeper
Galois filtration subgroup is used.
\end{proof}

\begin{corollary}[K17-NINE-DEGREES-PERIOD]
\label{cor:nine-degrees-period}
For every odd prime $\ell$ and every positive $k$ whose residue
modulo $(\ell-1)/2$ has a representative in
\[
                            \{0,1,2,3,4,5,6,7,8,9\},
\]
one has $K_{4k}(\mathbf Z)\{\ell\}=0$.
\end{corollary}

\begin{proof}
For a representative $j$ with $1\leq j\leq9$, (K32.2) and (K2)
give $A_\ell^{[-2j]}=0$. The equality
$k\equiv j\pmod{(\ell-1)/2}$ gives
$-2k\equiv-2j\pmod{\ell-1}$, so the two class-group components
are identical; (K2) gives the required vanishing. For $j=0$, the
component is trivial and vanishes by the extension-and-norm
argument following (K3). This proves only vanishing in these
residue classes, not periodicity of the orders of all $K$-groups.
\end{proof}

\paragraph{Arithmetic scope.}
These are seven new specified vanishing degrees beyond the known
degrees four and eight, conditional only on the stipulated Alp\"oge
arithmetic theorem and the established cited comparison theorems.
The finite range does not prove Vandiver at every prime and does
not remove the remaining all-prime tripod-kernel gate. The unitary
and modular-form calculations are complete exact certificates
with independent replays; they are not claimed to be Lean proofs.
